% Claim statement file for row claim_3_7 -- "TwoDisjointNode".
% Auto-generated stub by scaffold/solve_chapter.py at row start. The
% manager agent (or a worker it dispatches) fills the body below with the
% claim's statement(s) from the lecture notes -- statement only, no proof
% (proofs live in the sibling `claim_3_7_proof_*.tex` file).
%
% This file is a sibling of `main.tex` inside `<subsection>/tex/`. The
% `subfiles` package lets it render both standalone and as part of
% `main.tex`. Note: `main.tex` does NOT \subfile this statement file -- the
% sibling `_proof_` file restates the statement above the proof, so
% including both in `main.tex` would be redundant. This file exists for
% standalone reading / grep convenience.

\documentclass[main]{subfiles}

\begin{document}
% ref: claim_3_7 (statement)
% title: TwoDisjointNode
\def\rowref{claim\_3\_7}
\phantomsection\label{claim_3_7}

\begin{Lem}[Two disjoint node-splittings commute]
   Let $G=(J,V,E,L)$ be a CDMG and $W_1, W_2 \ins V$ two disjoint subsets of the output nodes of $G$.
   Then the CDMG obtained from first node-splitting $W_1$ and then node-splitting $W_2$ is the same CADMG that arises from first node-splitting $W_2$ and then node-splitting $W_1$:
   \[ \lp G_{\spl(W_1)} \rp_{\spl(W_2)} =  \lp G_{\spl(W_2)} \rp_{\spl(W_1)}  =  G_{\spl(W_1 \dcup W_2)}.   \]
\end{Lem}

\end{document}
